\section{Setup}
\label{sec:setup}

We describe the task, the selector, and the protocol. All code, configurations and per-run records are
released~(Appendix~\ref{app:repro}).

\subsection{Task: multi-query associative recall}

MQAR~\cite{arora2024zoology} presents a sequence of key--value pairs followed by queries; each query is a key seen
earlier, and the model must output its value. A model that answers correctly must retrieve a specific earlier token
pair by content, which is the ability that fixed-size recurrent states lose and attention keeps~\cite{arora2024zoology,
jelassi2024repeat}. We use the Zoology implementation (pinned commit \texttt{1ad20d1}) at sequence length
$L{=}2048$ with 128 key--value pairs and a vocabulary of 8{,}192. In this layout all 128 pairs occupy the first 256
tokens and the queries follow. We measure \emph{recall}: accuracy over query positions. Chance is about
$10^{-4}$.

\subsection{The selector}

The selector is a two-level version of the compression--selection--window design of NSA~\cite{yuan2025nsa}. Keys and
values are mean-pooled into blocks of $c{=}16$ tokens ($128$ blocks at $L{=}2048$), and blocks are grouped into
super-blocks of $g{=}11$ blocks ($12$ super-blocks). For each query and head:
\begin{enumerate}
\item score the query against every super-block summary that lies fully in the past, and keep the top $S_1$;
\item score the query against the $S_1 g$ block summaries inside the kept super-blocks; a softmax over these scores
forms a \emph{coarse} branch;
\item keep the top $K{=}4$ blocks and attend exactly over their $Kc{=}64$ raw tokens, the \emph{fine} branch;
\item attend over a causal sliding window of 32 tokens, the \emph{window} branch.
\end{enumerate}
Per-head sigmoid gates combine the three branches. The fine branch also adds the log-sigmoid of each selected block's
score to that block's token logits, which gives block scores a gradient path through the fine branch. The coarse
branch could in principle answer a query from mean-pooled summaries alone; Section~\ref{sec:selection} tests this. Selection uses hard top-$k$, so the super-block scores receive no
gradient through their indices; the shared query and key projections learn through the other branches. Invalid
(future or padded) candidates are masked with a large negative constant ($-10^9$). An independent code review noted
that this can misclassify a valid score at extreme magnitudes; telemetry over every M5-profile selector run (720
epochs of training forwards) recorded no valid score near the constant and no non-finite value, so the failure did
not occur there. The M4-profile runs recorded no such telemetry.

Table~\ref{tab:cost} compares the work of choosing what to read. Selector work per query is
$n_{\text{sup}} + \min(S_1, n_{\text{sup}})\,g$ scores, where $n_{\text{sup}}$ is the number of super-blocks that
lie fully in the query's past (12 at the end of a length-2048 sequence). This grows as $O(\sqrt{L})$ for fixed $S_1$ when
$g \approx \sqrt{L/c}$, giving $O(L^{1.5})$ per sequence, the same exponent as content-routed
attention~\cite{roy2021routing}. We report this count because it is the quantity the budget controls. \textbf{It
is not runtime}: our implementation has no fused kernel, and every selector setting ran 1.5--2.3$\times$
\emph{slower} per training step than dense attention.

\begin{table}[t]
\centering
\small
\begin{tabular}{lccc}
\toprule
Layer & selector scores per query & at $L{=}2048$ & seconds / step (L40S) \\
\midrule
Dense attention & -- (reads all $\le L$ past tokens) & -- & 0.0458 \\
Flat block selection (NSA-style) & $L/c$ & 128 & -- \\
Two-level selector, $S_1{=}2$ & $n_{\text{sup}} + S_1 g$ & 34 & 0.0668 \\
Two-level selector, $S_1{=}4$ & & 56 & 0.0744 \\
Two-level selector, $S_1{=}6$ & & 78 & 0.0820 \\
Two-level selector, $S_1{=}12$ (no pruning) & $n_{\text{sup}} + n_{\text{sup}} g$ & 144 & 0.1068 \\
\bottomrule
\end{tabular}
\caption{\textbf{The selector scores fewer candidates than flat block selection at small $S_1$, but runs slower
than dense attention in our implementation.} Scores counted per query and head, including masked candidates;
pooling, projections and attention over the selected tokens are excluded. Step times are from one smoke test per
setting on one NVIDIA L40S (M5 profile), without fused kernels.}
\label{tab:cost}
\end{table}

\subsection{Models and training recipe}

Every model has 2 layers, width 128 and 4 heads; only the sequence mixer changes. The dense baseline is Zoology's
multi-head attention, which applies dropout 0.1 inside attention; the selector has none. We train with AdamW
(learning rate $3.2\times10^{-4}$, weight decay 0.1, cosine schedule over 60 epochs, effective batch 16, no gradient
clipping) on 20{,}000 training sequences, and stop early once validation recall exceeds 0.98. The learning rate was
the only one, among those tried in earlier work in this project, that converged for attention models at
$L \ge 512$. We did not tune it per model.

\subsection{Protocol}

\paragraph{Metric.} Each run is scored on 2{,}000 held-out sequences that are never used for training or checkpoint
choice. The primary metric is held-out recall at the epoch of best validation recall (the earliest such epoch on a
tie); final-epoch recall is kept in every run record. MQAR is learned through an
abrupt phase change; we call a run \emph{transitioned} once its validation recall exceeds 0.30, and a run that never
does so within 60 epochs a \emph{non-transition}.

\paragraph{Two numerical profiles, never pooled.} The M4 experiments (Sections~\ref{sec:selection} and~\ref{sec:fixes}) used float32 and Zoology's original data generator, whose noise tokens draw from the global
random state (\emph{legacy} data). The M5 experiments (Sections~\ref{sec:budget} and~\ref{sec:dense}) used TF32 matrix
multiplication and a data generator with its own random stream (\emph{isolated} data), so that a seed fixes the data
exactly. Under the M5 profile, all models on a given seed trained and were tested on identical splits, which we verify
with SHA-256 hashes stored in each run's manifest.

\paragraph{Seeds, pre-registration and deviations.} Seed 123 is a pilot seed whose held-out split has informed design decisions
since M4; we label its results exploratory. Seeds 231 and 312 are further pilot seeds. Seeds 810001--810003 were
reserved on 1 October 2026 and not used before the M5 confirmation runs. Every training experiment's run list, metric and
decision rule were committed to the project's private development repository before its data existed. The public
release includes the dated protocol files but not that history, and such timestamps are self-attested. The selection diagnostic of Section~\ref{sec:selection} is post hoc. We
report every run, including non-transitions. Two deviations from the protocols occurred: the dense diagnostic runs on
seed 123 moved from one cloud provider to another (same GPU type and code) because no capacity was available, and we
stopped one float32 dense run (seed 810001) after 55 of its 60 pre-registered epochs, so it has no held-out score.
